\documentclass[../../main.tex]{subfiles}

\begin{document}

\section{不可测集}\label{section:ch02-不可测集}

我们已定义可测集且研究了可测集族的性质. 自然要问, 事实上, 是否存在不可测集? 这个问题的答案一点都不明显.

我们知道若集合 $E$ 的外测度为零, 则它是可测的, 且由于 $E$ 的任何子集外测度也为零, $E$ 的每个子集是可测的. 这是关于从集合的包含关系继承可测性所能得到的最好结果. 我们现在证明: 若 $E$ 是任何具有正外测度的实数集, 则 $E$ 存在不可测子集.

\begin{lemma}\label{lemma:ch02-平移零测}
令 $E$ 为有界的可测实数集. 假定存在有界、可数无穷的实数集 $\Lambda$ 使得 $E$ 的平移族 $\{\lambda+E\}_{\lambda\in\Lambda}$ 是不交的. 则 $m(E)=0$.
\end{lemma}

\begin{proof}
可测集的平移是可测的. 因此, 根据测度在可数不交并的可测集上的可数可加性,
\begin{align}
m\left[\bigcup_{\lambda\in\Lambda}(\lambda+E)\right]=\sum_{\lambda\in\Lambda}m(\lambda+E).\label{align:ch02-平移分解}
\end{align}
由于 $E$ 和 $\Lambda$ 都是有界集, 集合 $\bigcup_{\lambda\in\Lambda}(\lambda+E)$ 也是有界的, 因此有有限的测度. 于是\eqref{align:ch02-平移分解}的左边是有限的. 然而, 由于测度是平移不变的, 对每个 $\lambda\in\Lambda$, $m(\lambda+E)=m(E)\geqslant 0$. 因此, 由于集合 $\Lambda$ 是可数无穷的, 而\eqref{align:ch02-平移分解}右边的和式是有限的, 必有 $m(E)=0$.
\end{proof}

对于任何非空的实数集 $E$, 我们定义 $E$ 中的两个点为有理等价, 若它们的差属于有理数集 $\mathbf{Q}$. 容易看到这是一个等价关系, 即它是自反的、对称的与传递的. 我们称它为 $E$ 上的有理等价关系. 对于该关系, 存在 $E$ 的等价类族的不交分解. 谈到 $E$ 上的有理等价关系的选择集, 我们意味着集合 $\mathcal{C}_{E}$ 恰好由每个等价类的一个成员组成. 我们从选择公理推出存在这样的选择集. 选择集 $\mathcal{C}_{E}$ 由以下两个性质刻画:
\begin{enumerate}[(i)]
\item $\mathcal{C}_{E}$ 中的两个点的差不是有理的;
\item 对 $E$ 中的每个点 $x$, 存在 $\mathcal{C}_{E}$ 中的点 $c$ 使得 $x=c+q$, 其中 $q$ 是有理数.
\end{enumerate}

$\mathcal{C}_{E}$ 的第一个特征性质可方便地重述为:
\begin{align}
\text{对任意集合 }\Lambda\subseteq\mathbf{Q},\quad\{\lambda+\mathcal{C}_{E}\}_{\lambda\in\Lambda}\text{ 不交}.\label{align:ch02-选择集不交}
\end{align}

\begin{theorem}[Vitali]\label{theorem:ch02-Vitali}
任何具有正外测度的实数集 $E$ 包含一个不可测子集.
\end{theorem}

\begin{proof}
令 $E$ 为具有正的外测度的实数集. 根据外测度的可数次可加性, 我们可以假定 $E$ 是有界的. 令 $\mathcal{C}_{E}$ 为 $E$ 上的有理等价关系的任意选择集. 我们宣称 $\mathcal{C}_{E}$ 是不可测的. 为证明这个断言, 我们假设它是可测的且导出矛盾.

有理数可数无穷且在 $\mathbf{R}$ 中稠密. 令 $\Lambda_{0}$ 为任意有界可数无穷的有理数集. 由于 $\mathcal{C}_{E}$ 是可测的, 根据\eqref{align:ch02-选择集不交}, $\mathcal{C}_{E}$ 的通过 $\Lambda_{0}$ 的成员的平移族是不交的, 从\cref{lemma:ch02-平移零测}得出 $m(\mathcal{C}_{E})=0$. 然而, 由于 $\mathcal{C}_{E}$ 是 $E$ 上的有理等价关系的选择集, 我们有被包含关系
\begin{align}
E\subseteq\bigcup_{\lambda\in\mathbf{Q}}(\lambda+\mathcal{C}_{E}).\label{align:ch02-E覆盖}
\end{align}
从 Lebesgue 外测度的可数单调性、$\mathbf{Q}$ 的可数性以及测度的平移不变性推出
\begin{align*}
m^{*}(E)\leqslant\sum_{\lambda\in\mathbf{Q}}m(\lambda+\mathcal{C}_{E})=0.
\end{align*}
这与 $E$ 有正的外测度的假设矛盾, 从而定理得证.
\end{proof}

\begin{theorem}\label{theorem:ch02-不可加}
存在不交的实数集 $A$ 和 $B$ 使得
\begin{align*}
m^{*}(A\cup B)<m^{*}(A)+m^{*}(B).
\end{align*}
\end{theorem}

\begin{proof}
我们用反证法证明. 假设对每对不交的集合 $A$ 和 $B$, $m^{*}(A\cup B)=m^{*}(A)+m^{*}(B)$. 则根据可测集的最初定义, 每个集必须是可测的. 这与前一个定理矛盾.
\end{proof}

\begin{exercise}\label{exercise:ch02-29}
(i) 证明有理等价定义了任何集上的等价关系.

(ii) 对 $\mathbf{Q}$ 上的有理等价关系直接找到选择集.

(iii) 定义两个数为无理等价, 若它们的差是无理数. 它是 $\mathbf{R}$ 上的等价关系吗? 它是 $\mathbf{Q}$ 上的等价关系吗?
\end{exercise}

\begin{exercise}\label{exercise:ch02-30}
证明正外测度集合上的有理等价关系的任何选择集必须是不可数无穷的.
\end{exercise}

\begin{exercise}\label{exercise:ch02-31}
说明在 Vitali 定理的证明中仅须考虑 $E$ 是有界的情形就足够了.
\end{exercise}

\begin{exercise}\label{exercise:ch02-32}
若允许 $\Lambda$ 为有限或不可数无穷, \cref{lemma:ch02-平移零测} 仍然成立吗? 若允许 $\Lambda$ 为无界呢?
\end{exercise}

\begin{exercise}\label{exercise:ch02-33}
令 $E$ 为有限外测度的不可测集. 证明存在包含 $E$ 的 $G_{\delta}$ 集 $G$ 使得
\begin{align*}
m^{*}(E)=m^{*}(G),\qquad\text{而}\quad m^{*}(G\sim E)>0.
\end{align*}
\end{exercise}

\end{document}